\chapter{Conclusion}

The Academic Writing classes are over now. I would like to divide this conclusion into two movements. First, I will take a critical, broad view of the classes, specifically reflecting on the content, the professor's method, and whether my expectations were met. Then proceeding with the self-evaluation.

Looking back at my opening text, I can say that I learned a few things about writing and the academic world. Especially in the beginning, reading Tomitch's book on Academic Writing taught me a lot about how to think about language in general, reading, writing, what texts are, how to write summaries, and so on. We also checked the ``hamburger model'' for paragraph writing, went through many transition words and phrases, and learned about abstracts and their moves. We constantly received tips about writing, academia, and the English language. There is one point that was not clear in the expectations text I wrote: I was expecting more content. So much time was spent reading children's books, discussing topics unrelated (or only loosely related) to academic writing, and enduring long waiting times for the class to start. This happened both at the beginning of the class and after the break due to waiting until all students arrived, ultimately punishing those who got there on time. In general, productivity during class was unfortunately subpar.

The professor is highly qualified, has tons of experience in the academic world, and understands a lot about education and the English language. This was clear when we covered Tomitch and learned so many new and complex topics. Following the conjunction studies, the class slowed down significantly. Doing some simple calculations, the class is supposed to start at 6:30 pm and end around 9:40 pm (although officially it goes until 10:00 pm, most professors end a bit early due to students getting tired and some needing to catch the bus). The actual content starts on average at 7:30 pm, after waiting for everyone to get into class, waiting for the classmates to stop talking, and reading the children's book. The class goes until around 8:10 pm, and the professor asks for the class to return in 15 minutes. Usually, it takes around 30 minutes, so the content starts again at 8:40 pm. By 9:40 pm, the class already starts to leave. In total, we have around 80 minutes of actually productive class time. If the class started at 6:40 pm, went until 8:00 pm, returned at 8:15 pm, and then went until 9:30 pm, we would have 155 minutes of productive class time. Class time would almost double.

The syllabus includes: ``Reading and interpretation of academic texts in English; Scientific language; Structural and linguistic characteristics of the abstract genre; Structural and linguistic characteristics of the scientific article genre; Production of paragraphs, summaries, and abstracts''. We did in fact talk about everything, doing so briefly and superficially. I also noticed that the syllabus uses the plural form for ``paragraphs'', ``summaries'', and ``abstracts''. In practice, we only wrote one of each. Although technically a summary is a paragraph and an abstract is a type of summary, we didn't focus on them uniquely, making the plural in ``abstracts'' feel inadequate. Considering the available time, there was ample opportunity to read and study the whole Tomitch book, as well as to write texts applying every single conjunction, transition word, and phrase.

For me, the elephant in the room during class was the reading sessions at the beginning. I have a lot of difficulty understanding the reason for doing that. I understand that Professor Bailer also teaches English for Kids and that most of the class is composed of teachers or future teachers. That fact does not justify spending so much time every class on the same literary genre. As I discussed in the daily reflections, the books themselves are of great quality and bring valid points, alongside stories that are very simplistic. In an academic setting like ours, they feel shallow, the reflections consistently end up being too generic, and the analogies to academic writing are a bit vague. As I mentioned in my early reflections, there could be more variety. Incorporating poetry, philosophical and sociological texts (single chapters or paragraphs from authors such as Kant, St. Thomas Aquinas, Aristotle, and Nietzsche would go a long way), and scientific texts themselves would be highly beneficial. Trying to follow and understand tricky academic language from many different areas as a team during class would work as a highly productive icebreaker.

Another alternative is to use the beginning of the class for students to write texts with pen and paper right there, receiving immediate feedback from the professor. By this, I mean texts such as summaries, paragraphs, abstracts, and maybe even going beyond to write imaginary introductions, conclusions, and so on. I missed the writing part of the class a lot. The professor has immense knowledge and is excellent at correcting texts; we only made use of that skill during the three writing activities we had. Allocating a small amount of time every class to write something alongside the story time, followed by live feedback from the professor, would help us improve much more. And I mainly suggest to keep the exercises and writing in the class itself, not in homework form. That way the students can keep more focused, will not cheat and can use their free time at home for other activities.

Perhaps a good solution would be to mix and match different types of texts, including short writing sessions alongside the reading ones. In that scenario, we could get familiar with different ways of thinking and writing. This would allow us to go beyond the self-help conclusions of the books, like ``be yourself'', ``be creative'', and ``love others'', to think at least a little about ``what it means to \textbf{be}'', ``what the `I' is'', ``how to be creative while keeping a foot in reality'', ``what the definition of love is'', and so on.

Now, regarding the self-reflection part, I will continue the texts by trying to answer the questions presented on the last day of class. I can indeed follow the structure of abstracts, navigate the sections of academic texts, deal with scientific language, write paragraphs, summaries, and abstracts, and make corrections based on feedback. However, I do not feel particularly more confident than before. I already had a little experience with academic writing from high school (I even published an article that can be found here: \url{https://scholar.google.com.br/scholar?hl=pt-BR&as_sdt=0%2C5&q=two+kingdoms+gustavo+guerreiro&btnG=}). Therefore, I already knew much of what we saw. Also, in my Computer Science major, I had already taken the ``Produção Textual Acadêmica'' course. That course was online and contained texts with only 15 minutes of reading time, along with four in-person meetings with Professor Waldir Luiz Hellmann. Those encounters were productive, though; we had many exercises to complete and discuss together, alongside numerous discussions about the topics of the academic articles themselves. Those four meetings made me think that this class would contain 16 similar sessions. It is only natural that different professors will have different methods, and perhaps having hyper-productive classes every week could be really tiring.

Regarding reading, I don't think I got better or worse; in English, the academic language isn't particularly a problem for me. Maybe revisiting the idea of using the abstract to evaluate if the paper will be useful to me before reading the rest is a good tip, although I usually already did that automatically. My main difficulty lies in searching for a good paper, rather than the reading process itself. Probably, the real challenge is being able to skim over the abstract as fast as possible in order to understand if I want to continue working with that paper or move on to the next.

Regarding writing, I cannot say that this class improved it much; we had way too little written production for me to actually be able to compare how I wrote before the class began and after. I believe that I still need to keep working on my conjunction use and get rid of my comma splices. Those are two of the points with which I struggle the most. The assignment that best represents my learning has to be the abstract we wrote. After all, it uses all of what we learned: it is a paragraph, it is a summary, and it uses clear indicators of the moves and conjunctions.

Regarding the language learning aspect, my use of conjunctions, punctuation, and general formal understanding got a bit better. I felt that I was already quite familiar with the language, to the point where I wouldn't say I struggle with English itself. My struggles tend to lean towards writing in general. The habits that contributed the most were definitely practicing, reading, writing, listening, and speaking; slowly, the language gets more solid.

At the beginning of this semester, I would have told myself to focus more on the content of the class and note all the most important points discussed since day one. I felt that at the start I was a bit lost regarding the concept of how to discuss or work on the Reflective Portfolio, so I ended up not taking many notes or thinking too much about the content.

I plan to continue improving my academic writing in English by reading more, writing more, reviewing what we saw, and finishing Tomitch's book so I can produce excellent academic work during my upcoming studies abroad.

Finally, I would not like the conclusion to end on a negative note. Instead, I would like to reflect on this last semester and how Academic Writing was present in it. This was a very atypical semester. I had already finished all the undergraduate studies demanded by my Computer Science major and was just preparing for my classes abroad, which led me to study subjects such as Civil Law and Academic Writing. Despite my previous criticism and self-assessment in learning, I cannot say that I disliked the class. The classes were calmer; I knew what was expected of me, and I found the experience of writing the portfolio quite enjoyable. I guess that is why I wrote so much. I never felt discomfort or stress caused by Academic Writing. Professor Bailer also showed love for what she does; I feel genuine care and interest in the students coming from her. She cares about the quality of the class and the well-being of the students, and she always wants to bring good materials and provide a useful and enjoyable experience with ethics and enthusiasm, independent of the circumstances. Surely, this was not the final class I thought I would take when I enrolled in Computer Science back in February 2022, but it was a good one. Anyway, I am not sure if my observations given in this conclusion are shared by the majority of my colleagues or not; either way, I hope that at least some of my comments and observations make sense and help to improve the classes even further. 
\\ \\
\textit{Thank you for your time in reading this portfolio and for your dedication throughout my final semester.}